\chapter{Limit Theorem}

\section{The Law of Large Numbers}

\begin{theorem}[Bernoulli's Law of Large Numbers]
  Suppose that \(n_A\) is the number of successes in \(n\) Bernoulli trials with success probability \(p\). Then for any \(\varepsilon > 0\), we have
  \[
    \lim_{n \to \infty} P\!\left\{ \left| \frac{n_A}{n} - p \right| < \varepsilon \right\} = 1.
  \]
\end{theorem}

\begin{theorem}[Chebyshev's Law of Large Numbers]
  Let \(X_1, X_2, \ldots\) be an independent sequence with finite variance, and there exists a constant \(K\) such that \(\Var(X_i) \leq K\) for all \(i\). Then for any \(\varepsilon > 0\), we have
  \[
    \lim_{n \to \infty} P\!\left\{ \left| \frac{1}{n} \sum_{i=1}^n X_i - \frac{1}{n} \sum_{i=1}^n \E(X_i) \right| < \varepsilon \right\} = 1.
  \]
\end{theorem}

\begin{theorem}[Khinchin's Law of Large Numbers]
  Let \(X_1, X_2, \ldots\) be an independent and identically distributed (IID) sequence with \(\E(X_i^k) = \mu_k < \infty\) for all \(k \in \N\). Then for any \(\varepsilon > 0\), we have
  \[
    \lim_{n \to \infty} P\!\left\{ \left| \frac{1}{n} \sum_{i=1}^n X_i^k - \mu_k \right| < \varepsilon \right\} = 1.
  \]
\end{theorem}

\begin{example}
  Let \(X_1, X_2, \ldots\) be an IID sequence with \(\E(X_i) = \mu\) and \(\Var(X_i) = \sigma^2\). Then for any \(\varepsilon > 0\), we have
  \[
    \frac{1}{n} \sum_{i=1}^n X_i^2 \xrightarrow{P} \sigma^2 + \mu^2.
  \]
\end{example}

\section{The Central Limit Theorem}

\begin{theorem}[Levy's Central Limit Theorem]
  Let \(X_1, X_2, \ldots\) be an IID sequence with \(\E(X_i) = \mu\) and \(\Var(X_i) = \sigma^2\). Then for any \(x \in \R\), we have
  \[
    \lim_{n \to \infty} P\!\left\{ \frac{\sum_{i=1}^n X_i - n\mu}{\sigma \sqrt{n}} \leq x \right\} = \Phi(x),
  \]
  where \(\Phi(x)\) is the c.d.f.\@ of the standard normal distribution. Equivalently, we have
  \[
    \frac{\sum_{i=1}^n X_i - n\mu}{\sigma \sqrt{n}} \sim \mathrm{N}(0, 1), \quad \sum_{i=1}^n X_i \sim \mathrm{N}(n\mu, n\sigma^2).
  \]
\end{theorem}

\begin{example}
  In order to protect the sovereignty of territorial waters, we prepare to launch 900 missile strikes against enemy targets. Suppose that the number of missiles hitting the target \(X_k\) in \(k\)-th strike follows identical distribution and is independent of each other. If \(\E(X_k) = 4\) and \(\Var(X_k) = 1.21\) \((k = 1, 2, \dots, 900)\), then calculate the probability of the event that the total number of missiles hitting the target is between 3534 and 3699.
\end{example}
\begin{solution}
  By using Levy's Central Limit Theorem, we have
  \[
    \sum_{i=1}^{900} X_i \sim \mathrm{N}(900 \times 4, 900 \times 1.21) = \mathrm{N}(3600, 1089).
  \]
  Then we can calculate the probability as follows:
  \begin{align*}
    P\!\left\{ 3534 \leq \sum_{i=1}^{900} X_i \leq 3699 \right\} &\approx \Phi\!\left( \frac{3699 - 3600}{\sqrt{1089}} \right) - \Phi\!\left( \frac{3534 - 3600}{\sqrt{1089}} \right) \\
    &= \Phi(3) - \Phi(-2) = \Phi(3) + \Phi(2) - 1. \qedhere
  \end{align*}
\end{solution}

\begin{theorem}[De Moivre-Laplace Central Limit Theorem]
  Suppose that \(Y_n \sim \mathrm{B}(n, p) (0 < p < 1)\). Then for any \(x \in \R\), we have
  \[
    \lim_{n \to \infty} P\!\left\{ \frac{Y_n - np}{\sqrt{np(1-p)}} \leq x \right\} = \Phi(x).
  \]
\end{theorem}

\begin{theorem}[Liapunov's Central Limit Theorem]
  Let \(X_1, X_2, \ldots\) be an independent sequence with \(\E(X_i) = \mu_i\) and \(\Var(X_i) = \sigma_i^2 \neq 0\) for all \(i\). If there exists a constant \(\delta > 0\) such that
  \[
    \lim_{n \to \infty} \frac{\sum_{i=1}^n \E\!\left( |X_i - \mu_i|^{2+\delta} \right)}{\left( \sum_{i=1}^n \sigma_i^2 \right)^{1 + \delta/2}} = 0,
  \]
  then for any \(x \in \R\), we have
  \[
    \lim_{n \to \infty} P\!\left\{ \frac{\sum_{i=1}^n X_i - \sum_{i=1}^n \mu_i}{\sqrt{\sum_{i=1}^n \sigma_i^2}} \leq x \right\} = \Phi(x).
  \]
\end{theorem}